% --- Archon physics formalization source begin ---
% archon:physics
% archon:chemistry
% archon:covers IChO2026Problems/problem_icho_2026_t2_a2.lean
% archon:source-report reports/icho_2026/problem_icho_2026_t2_a2.source.json
% archon:problem-id icho_2026_t2
% archon:part-id T2-A2

\chapter{Icho Chemistry problem icho\_2026\_t2\_a2}
\label{ch:IChO2026Problems_problem_icho_2026_t2_a2}

\paragraph{Problem source.}
\#\# Source
58th International Chemistry Olympiad, Tashkent, Uzbekistan, 2026.
Official-materials index: https://www.icho2026.uz/problems
English problem PDF: ../icho\_2026\_source/raw/theory\_problem.pdf
English solution/rubric PDF: ../icho\_2026\_source/raw/theory\_solution.pdf
Problem PDF page 16; printed local question page 2; solution starts on PDF page 16.
The page PNG is primary visual evidence for formulas, structures, tables, and figures.

\#\# T2. Kinetics of the Belousov-Zhabotinsky Reaction
T2. Kinetics of the Belousov-Zhabotinsky reaction 6\% When the chemicals KBrO3, CH2(COOH)2 (malonic acid, – MA), Ce(SO4)2, and H2SO4 are mixed, the concentration, 𝐶, of Ce species and the colour of the solution oscillates between yellow (due to Ce4+) and colourless (due to Ce3+). A schematic picture is given below: 2.1 Write the overall balanced reaction equation of the Belousov-Zhabotinsky reaction (BZ reaction). Note: in this reaction, MA is oxidised to CO2, BrO− 3 is reduced to Br−, and Ce(IV) is a catalyst. The mechanism of the BZ reaction is shown below (BMA −CHBr(COOH)2): Process A ⎧ ⎪ ⎪ ⎨ ⎪ ⎪ ⎩ (1) HBrO2 + BrO− 3 + H+ 𝑘1 −−→2BrO⋅ 2 + H2O (2) BrO⋅ 2 + Ce3+ + H+ 𝑘2 −−→HBrO2 + Ce4+ (3) 2HBrO2 𝑘3 −−→BrO− 3 + HBrO + H+ 𝑘1 = 1.0 × 104 M−2 s−1 𝑘2 = 6.2 × 104 M−2 s−1 𝑘3 = 4.0 × 107 M−1 s−1 Process B ⎧ ⎪ ⎪ ⎨ ⎪ ⎪ ⎩ (4) HBrO2 + Br−+ H+ 𝑘4 −−→2HBrO (5) BrO− 3 + Br−+ 2H+ 𝑘5 −−→HBrO + HBrO2 (6) HBrO + MA 𝑘6 −−→BMA + H2O 𝑘4 = 2.0 × 109 M−2 s−1 𝑘5 = 2.1 M−3 s−1 𝑘6 = 8.2 M−1 s−1 Process C \{(7) Ce4+ + BMA 𝑘7 −−→Ce3+ + Br−+ other products 𝑘7 = 1.0 × 102 M−1 s−1 In the BZ reaction, Process C occurs continuously, while Processes A and B alternate with each other. Assume that: - when Process A occurs the solution is yellow and Process B practically does not occur; - when Process B occurs the solution is colourless and Process A practically does not occur; - [BrO3−]0 = 0.06 M, [MA]0 = 0.1 M, [H+]0 = 0.8 M, [Ce4+]0 = 0.001 M; - the concentrations of the reactants (except Ce4+) and pH are maintained constant throughout the BZ reaction.

\#\# Current subquestion 2.2
Using the steady-state approximation, calculate the stationary molar concentrations of HBrO2 in Processes A and B, [HBrO2]A and [HBrO2]B.

\#\# Dataset metadata
IChO 2026; T2; raw rubric points=11; kind=theory; formalization\_ready=true.

\paragraph{Shared problem context.}
T2. Kinetics of the Belousov-Zhabotinsky reaction 6\% When the chemicals KBrO3, CH2(COOH)2 (malonic acid, – MA), Ce(SO4)2, and H2SO4 are mixed, the concentration, 𝐶, of Ce species and the colour of the solution oscillates between yellow (due to Ce4+) and colourless (due to Ce3+). A schematic picture is given below: 2.1 Write the overall balanced reaction equation of the Belousov-Zhabotinsky reaction (BZ reaction). Note: in this reaction, MA is oxidised to CO2, BrO− 3 is reduced to Br−, and Ce(IV) is a catalyst. The mechanism of the BZ reaction is shown below (BMA −CHBr(COOH)2): Process A ⎧ ⎪ ⎪ ⎨ ⎪ ⎪ ⎩ (1) HBrO2 + BrO− 3 + H+ 𝑘1 −−→2BrO⋅ 2 + H2O (2) BrO⋅ 2 + Ce3+ + H+ 𝑘2 −−→HBrO2 + Ce4+ (3) 2HBrO2 𝑘3 −−→BrO− 3 + HBrO + H+ 𝑘1 = 1.0 × 104 M−2 s−1 𝑘2 = 6.2 × 104 M−2 s−1 𝑘3 = 4.0 × 107 M−1 s−1 Process B ⎧ ⎪ ⎪ ⎨ ⎪ ⎪ ⎩ (4) HBrO2 + Br−+ H+ 𝑘4 −−→2HBrO (5) BrO− 3 + Br−+ 2H+ 𝑘5 −−→HBrO + HBrO2 (6) HBrO + MA 𝑘6 −−→BMA + H2O 𝑘4 = 2.0 × 109 M−2 s−1 𝑘5 = 2.1 M−3 s−1 𝑘6 = 8.2 M−1 s−1 Process C \{(7) Ce4+ + BMA 𝑘7 −−→Ce3+ + Br−+ other products 𝑘7 = 1.0 × 102 M−1 s−1 In the BZ reaction, Process C occurs continuously, while Processes A and B alternate with each other. Assume that: - when Process A occurs the solution is yellow and Process B practically does not occur; - when Process B occurs the solution is colourless and Process A practically does not occur; - [BrO3−]0 = 0.06 M, [MA]0 = 0.1 M, [H+]0 = 0.8 M, [Ce4+]0 = 0.001 M; - the concentrations of the reactants (except Ce4+) and pH are maintained constant throughout the BZ reaction.

\paragraph{Current subquestion.}
Using the steady-state approximation, calculate the stationary molar concentrations of HBrO2 in Processes A and B, [HBrO2]A and [HBrO2]B.

\paragraph{Recorded answer/context.}
𝑑[BrO⋅ 2] 𝑑𝑡 = 2𝑘1[HBrO2][BrO− 3 ][H+] −𝑘2[BrO⋅ 2][Ce3+][H+] = 0 (eq.2.1) 𝑑[HBrO2] 𝑑𝑡 = −𝑘1[HBrO2][BrO− 3 ][H+]+𝑘2[BrO⋅ 2][Ce3+][H+]−2𝑘3[HBrO2]2 = 0 (eq.2.2) When adding equations 2.1 and 2.2 we get the following: 𝑘1[HBrO2][BrO− 3 ][H+] −2𝑘3[HBrO2]2 = 0 (eq.2.3) [HBrO2]𝐴= 𝑘1 2𝑘3 [BrO− 3 ][H+] (eq.2.4) [HBrO2]𝐴= 1.0×104 2×4.0×107 × 0.06 × 0.8 = 6.0 × 10−6 M The steady-state approximation for Process B: 𝑑[HBrO2] 𝑑𝑡 = −𝑘4[HBrO2][Br−][H+] + 𝑘5[BrO− 3 ][Br−][H+]2 = 0 (eq.2.5) [HBrO2]𝐵= 𝑘5 𝑘4 [BrO− 3 ][H+] (eq.2.6) [HBrO2]𝐵= 2.1 2.0×109 × 0.06 × 0.8 = 5.04 × 10−11 M [HBrO2]𝐴= 6.0 × 10−6 M [HBrO2]𝐵= 5.04 × 10−11 M For fully correct equations 2.1, 2.2 and 2.5: 2 points each. If only one incorrect coefficient in front of the rate constant in equations 2.1, 2.2 and 2.5: 1 point each. For corresponding equations 2.3, 2.4 and 2.6: 1 point each. For calculated [HBrO2]A and [HBrO2]B values: 1 point each. For correct calculation without formula derivation: full marks. 11.0 pt

\paragraph{Figure/image paths.}
\begin{itemize}
\item ../icho\_2026\_source/image/T2\_page-2.png
\item ../icho\_2026\_source/image/T2\_page-1.png
\end{itemize}

\paragraph{Formalization target.}
create a compiling Lean file with sorry bodies at `IChO2026Problems/problem\_icho\_2026\_t2\_a2.lean`.
The Lean declarations must preserve every mathematical object, quantifier, hypothesis, side condition, approximation/error guarantee, and requested conclusion from the source statement.
Use Mathlib/Physlib/Chemistry names found through LeanExplore where available. Prefer the configured benchmark Base library; introduce a local definition only when the source genuinely requires an object that the environment does not expose.

\begin{theorem}[Icho Chemistry formalization target]
\label{thm:physics:icho_2026_t2_a2:target}
The assigned autoformalize agent should translate this IChO chemistry problem into Lean declarations in the covered file, with theorem and lemma proof bodies written as `by sorry`.
\end{theorem}
\begin{proof}
This is an autoformalization task, not a proof task. Produce faithful statements that can later be proved without weakening the source contract.
\end{proof}
% --- Archon physics formalization source end ---
